% ==============================================================
% Input Appendix configuration
% ==============================================================
% appendix.tex
% ==============================================================
\appendix

% --- Force letter-based section counter ---
\renewcommand{\thesection}{\Alph{section}}
\renewcommand{\thesubsection}{\Alph{section}.\arabic{subsection}}
\renewcommand{\thesubsubsection}{\Alph{section}.\arabic{subsection}.\arabic{subsubsection}}

% --- Reset section counter ---
\setcounter{section}{0}

% --- Restart and prefix all counters ---
\setcounter{equation}{0}
\setcounter{figure}{0}
\setcounter{table}{0}
\setcounter{lstlisting}{0}

\renewcommand{\theequation}{\Alph{section}.\arabic{equation}}
\renewcommand{\thefigure}{\Alph{section}.\arabic{figure}}
\renewcommand{\thetable}{\Alph{section}.\arabic{table}}
\renewcommand{\thelstlisting}{\Alph{section}.\arabic{lstlisting}}

% --- Opening title ---
\pagesectionvspace
\section*{\trAppendices}
\phantomsection
\addcontentsline{toc}{section}{\trAppendices}
\label{sec:appendix}
\vspace{10mm}

% --- Typography: smaller font for appendix content ---
\small

% adjust width
\begin{adjustwidth}{\appendixtab}{0mm}

% ==============================================================
% Appendix A
% ==============================================================
\section{Supplementary Data} \label{app:data}

This is appendix \Cref{app:data}. The \texttt{appendix} command
resets the section counter and switches labels to letters
automatically.

\lipsum[1]

\subsection{First subsection} \label{app:data:sub1}

This is \Cref{app:data:sub1}. Subsections follow the \texttt{A.1},
\texttt{A.2} pattern automatically.

\lipsum[1]

\subsection{Supplementary figure} \label{app:data:fig}

The \Cref{fig:app:example} is a supplementary figure. Note the
label follows the \texttt{A.1} numbering convention automatically.

\lipsum[1-3]

\begin{center}
\begin{minipage}{\linewidth}
\centering
\includegraphics[width=0.6\linewidth]{example-image-golden}
\captionof{figure}[Supplementary figure]
{
    \textbf{Supplementary figure.}
    Ut enim ad minim veniam, quis nostrud exercitation ullamco
    laboris nisi ut aliquip ex ea commodo consequat.
    \;\textbf{a}\;---\;First panel description.
    \;\textbf{b}\;---\;Second panel description.
}
\label{fig:app:example}
\end{minipage}
\end{center}

\lipsum[1-3]

% ==============================================================
% Appendix B
% ==============================================================
\section{Derivations} \label{app:derivations}

This is appendix \Cref{app:derivations}. Cross-references from the
main body work normally, for example back to \Cref{sec:equations}.

\lipsum[1]

\subsection{First derivation} \label{app:derivations:sub1}

\lipsum[1]

\begin{equation}
    \label{eq:app:example}
    f(x) = \int_{-\infty}^{\infty} \hat{f}(\xi)\, e^{2\pi i \xi x}\, d\xi
\end{equation}

\Cref{eq:app:example} is an equation inside an appendix --- it is
cited and numbered normally.

\subsection{Supplementary table} \label{app:derivations:tbl}

The \Cref{tbl:app:example} is a supplementary table. Columns
follow the same booktabs convention as the main body.

\lipsum[1]

\begin{center}
\begin{threeparttable}
\scriptsize
\sffamily
\captionof{table}[Supplementary parameter table]{
    Derived parameters referenced in \Cref{app:derivations}.
    Values computed analytically from \Cref{eq:app:example}
    using the method described in \Cref{app:derivations:sub1}.
}
\label{tbl:app:example}
\rowcolors{2}{white}{AccentColorVeryLight}
\begin{tabular}{
    >{\raggedright\arraybackslash}m{2.3cm}
    >{\raggedright\arraybackslash}m{2.3cm}
    >{\raggedright\arraybackslash}m{2.3cm}
    >{\raggedright\arraybackslash}m{2.3cm}
    >{\raggedright\arraybackslash}m{2.3cm}}
\toprule
\textbf{Parameter} & \textbf{Symbol} & \textbf{Value} & \textbf{Unit} & \textbf{Notes} \\
\midrule
Frequency       & $\xi$      & 1.234 & Hz           & Measured          \\
Amplitude       & $\hat{f}$  & 0.876 & ---          & Normalised        \\
Phase shift     & $\phi$     & 0.314 & rad          & See note\tnote{a} \\
Decay constant  & $\lambda$  & 2.718 & s$^{-1}$     & Derived           \\
Derived result  & $f(x)$     & 3.141 & ---          & See note\tnote{b} \\
\bottomrule
\end{tabular}
\begin{tablenotes}
    \item[\textit{a}] Phase shift computed relative to the reference
    signal defined in \Cref{eq:app:example}.
    \item[\textit{b}] Final result after applying the Fourier transform
    over the full integration domain $[-\infty, \infty]$.
\end{tablenotes}
\end{threeparttable}
\end{center}

\lipsum[1]

% ==============================================================
% Appendix C
% ==============================================================
\section{Code Listings} \label{app:listings}

\lipsum[1]

\begin{lstlisting}[caption={Example function}, label={lst:app:example}]
def compute(x: float) -> float:
    """Compute the square root safely."""
    if x < 0:
        raise ValueError(f"Expected non-negative, got {x}")
    return x ** 0.5
\end{lstlisting}

\Cref{lst:app:example} is a listing inside an appendix.

\lipsum[1]

\end{adjustwidth}